\section*{Conclusion}

A genomic resistance model can look excellent and still be measuring the wrong thing. Our own model
reached 0.850 average precision on a random split, led every external baseline, and was matched to
within 0.008 by one of its own components; a model that read no sequence at all reached 0.705 on the
same split. Holding out whole lineages dropped it to 0.645, of which two-fifths was the metric
responding to a lower resistance prevalence rather than the model losing skill, and the remaining
loss fell on four of seven drugs while three were untouched. On those four, very major error rates
rose above two thirds while categorical agreement stayed flat or improved --- the measure most
papers report moving in the opposite direction from the measure a laboratory would use.

The cause was the features, not the model class. Replacing a reference pan-genome catalogue with
each strain's own determinant calls recovered most of the loss on exactly the drugs whose
determinants are acquired, and left the point-mutation control unchanged. Once the determinants are
present, a sparse linear model on them is enough: it matched gradient boosting on unseen lineages,
beat a fixed determinant rule by 0.104 on 10,232 external isolates where gradient boosting managed
0.065, and was at or above the rule for every drug. In a second species whose determinants are
nearly sufficient on their own, no model beat the rule at all.

What we would ask of the next genomic resistance paper follows from that: a determinant rule as the
per-drug baseline, a metadata-and-assembly model reported on every split, a lineage-disjoint test
read with its prevalence shift removed, a simple learned baseline beside any complex one, and
clinical error rates alongside the ranking metrics. Our own architecture contributed a small,
seed-stable increment to a stack and no interpretable account of its predictions; the evaluation
protocol is what made the two-group finding measurable, and that finding --- that genomic
resistance prediction in \Ecoli{} is not one problem but two, with a sharp and drug-specific
boundary between them --- is what we would keep.

\section*{Data availability}
Code, pre-registration documents, Kaggle notebooks and result tables are available at
\url{https://github.com/FaiazPronoy/TriLens-AMR} \todo{archive a tagged release on Zenodo and give
its DOI}. The training data are the Arcadia Science release
\cite{arcadia2024dataset,bell2025arcadia}; the external cohort, its determinant calls and its MICs
are public in NCBI Pathogen Detection \cite{ncbi_pathogen_detection}; the Giessen collection is that
of Ren et al.\ \cite{ren2022giessen}. Exact table versions and file checksums are recorded in the
repository. Every split, filter and model is driven by a single fixed seed (42), so every figure
reported here is reproducible, with the training-variability caveat stated in the Limitations.

\section*{Supplementary data}
Supplementary Data are available at NARGAB Online: Supplementary Methods (the full TriLens
specification with every equation; the presence/absence matrix repair; the per-gene statistics and
lineage encoding), all pre-registration documents with their amendments, and per-drug tables for
every experiment.

\section*{Acknowledgements}
\todo{}

\section*{Author contributions}
\todo{CRediT roles.}

\section*{Funding}
\todo{}

\section*{Conflict of interest}
\todo{None declared, if true.}
